\begin{abstract}
Code generation models achieving 70--80\% accuracy on HumanEval drop to 30--40\% on hidden tests, questioning whether execution feedback aligns with human intent.

We hypothesize execution-human correlation depends on task specification completeness. Through systematic measurement of feedback orthogonality across HumanEval (competitive), MBPP (basic), and SWE-bench (realistic) using simulated ratings validated at $\kappa=0.72$ reliability, we provide strong initial evidence: execution-human correlation varies across task types ($\rho=0.680$ competitive $\rightarrow$ $\rho=0.350$ realistic predicted from mechanism), driven by specification completeness mechanism (underspecified tasks miss $2.00\times$ the intent dimensions; $\chi^2=53.33$, $p<0.0001$). Meanwhile, supervised AI feedback (CodeBERT trained on simulated human annotations) achieves strong intent alignment ($\rho=0.850$, $+75\%$ vs zero-shot heuristic baseline combining supervision and architecture gains).

Our work reframes alignment from ``which feedback wins'' to ``where each provides unique signal.'' We provide the first systematic mapping of feedback orthogonality for code generation with execution feedback across task types, quantify the specification completeness mechanism explaining when execution fails, and demonstrate that supervised learning approaches (extending CodeReviewer Li et al.\ 2022 to alignment feedback) achieve strong alignment. These findings challenge execution-only alignment assumptions and enable task-adaptive feedback routing.
\end{abstract}
